\section{Limitations}
\label{sec:limitations}

\paragraph{Simulated buyers are not human buyers.} VEB's committee is
LLM-simulated. This buys control, reproducibility, and matched pairs, but a
simulated buyer may be exploitable in ways a human is not, or may reward
rhetorical patterns a real executive would resist. We mitigate with
personality-grounded personas, private notes hidden from the seller, gated
disclosure that requires earned moves, and stochastic re-sampling; we do not claim
that VEB performance transfers one-to-one to live selling. Establishing
sim-to-real transfer is important future work.

\paragraph{Conflict of interest, stated plainly.} VEB is built by the authors of
the sales methodology it evaluates, and the methodology-pack track tests that same
methodology. A skeptical reader should treat the \emph{existence} of a pack lift as
a claim made by an interested party, and we have designed the benchmark so that the
claim does not have to be taken on trust: the grid, transcripts, scorecards, seeds,
and scoring code are all released; the lift is computed by deterministic,
inspectable math over evidence-cited scorecards; and the results themselves are not
flattering marketing---the pack \emph{hurts} some models (Table~\ref{tab:lift}
shows a negative mid-tier lift), helps only where behavioral markers move, and the
headline pooled effect is small. We report what the data shows, including where it
undercuts the sales pitch. Independent replication is the real remedy, and the
release is structured to make replication cheap: every number in this paper can be
recomputed from the published artifacts without contacting us.

\paragraph{Wins concentrate on few behaviors.} Our own analysis
(Section~\ref{sec:analysis}) shows that reaching the economic buyer and holding
price predict most clean wins. This is consistent with how the underlying
methodology says enterprise deals work, but it also means the leaderboard partially
measures whether a model happens to trigger the specific behaviors the simulator
rewards. We surface this rather than hide it: the behavioral decomposition
(Figure~\ref{fig:behavior}) lets a reader separate ``executes the rewarded
behaviors'' from ``scores well,'' and the failure-mode audit trail shows
\emph{which} behaviors gate which outcomes. Whether those gates match live buying
committees is exactly the sim-to-real question above.

\paragraph{Scenario admission is a selection effect on the pack lift.} Scenario
calibration (Section~\ref{sec:environment}) admits a scenario only if a
disciplined reference seller \emph{executing the methodology under test} can win
it. This guarantees every deal is winnable, but it also means the pack track is
evaluated exclusively on terrain the methodology is certified to clear: by
construction there can be no scenario where following the methodology is the
wrong strategy. The measured pack lift is therefore an estimate of ``does
injecting the methodology help on methodology-winnable deals,'' not ``does the
methodology beat alternatives on arbitrary deals.'' A symmetric design would
calibrate scenarios against several distinct selling philosophies; we flag this
as the natural extension for the multi-methodology track already implied by the
pack mechanism.

\paragraph{Rubric commitment.} The grading rubric encodes one particular
enterprise-sales methodology. A model tuned to a different but effective selling
philosophy could be under-credited. We regard this as a scoped design decision:
VEB measures competence \emph{against a specified, published standard}, and the
standard is transparent and auditable rather than implicit in an opaque reward.

\paragraph{Judge model and correlated error.} The automated judge is itself an
LLM and can share blind spots with the models it grades, especially same-family
models (self-preference). Our quote-or-abstain constraint, deterministic roll-up,
and the forthcoming human validation with a blind arbitration subset are designed to
bound this risk, though they do
not eliminate it; residual judge--model correlation is reported, not assumed away.
As a direct probe we re-grade the entire grid under a hardened judge that
reconstructs the objective dimensions deterministically from the event log
(Section~\ref{sec:hardened}); mean SQS moves only \resHardenedMeanDelta{}~points and
the top of the board is unchanged, but this bounds rather than eliminates the shared-
prior concern, since both judges are still LLMs on the subjective dimensions.

\paragraph{Validation sample size.} Human expert review is expensive, so
validated agreement is estimated on a stratified sample, not the whole grid.
Agreement statistics therefore carry sampling uncertainty, which we report with
confidence intervals; rare strata (e.g., hard-tier wins) have the widest
intervals.

\paragraph{Coverage.} \numScenarios{} scenarios across \numIndustries{} industries
is broad but finite, and each scenario, once public, can in principle be studied
and overfit. The stochastic buyer and per-seed re-sampling raise the cost of
overfitting but do not make the environment adversarially unforgeable; we treat
periodic scenario refresh as part of the benchmark's maintenance.

\paragraph{English, text-only, and single-seller.} The current release is
English-language, text-channel, and evaluates a single seller agent against the
committee. Voice, multilingual selling, and multi-seller team dynamics are out of
scope for this version.
